\begin{frame}[fragile]{실습: 미니 RAG --- 문서 벡터 만들기}
  14.2의 \texttt{train\_skipgram}·\texttt{cos\_sim}을 그대로 쓰고,
  문서 안 단어 벡터를 \textbf{평균}해 ``문서 벡터''를 만드는
  \texttt{mean\_vec}를 추가. \; 에너지·동물·컴퓨팅 5개 문서:
\begin{lstlisting}
# 문서 벡터 = 존재하는 단어 벡터의 평균
def mean_vec(words, emb):
    dim = len(next(iter(emb.values())))
    present = [w for w in words if w in emb]
    if not present:
        return None
    v = [0.0] * dim
    for w in present:
        for k in range(dim):
            v[k] += emb[w][k]
    return [x / len(present) for x in v]

tok = lambda s: s.replace(". ", " ").split()
corpus_docs = {
    "D1 전기차": "전기차는 배터리로 움직인다. ...",
    "D2 태양광": "태양광 발전은 햇빛으로 전기를 만든다. ...",
    "D3 양자":   "양자 컴퓨터는 큐비트로 정보를 처리한다. ...",
    "D4 고양이": "고양이는 애묘인이 좋아하는 동물이다. ...",
    "D5 배터리": "배터리와 태양광 발전은 ...",
}
flat = []
for _n, t in corpus_docs.items():
    flat += tok(t)
random.seed(42)
emb = train_skipgram(flat, window=2, dim=8, epochs=400, lr=0.1, neg_k=3)
docvecs = {n: mean_vec(tok(t), emb) for n, t in corpus_docs.items()}
# 토큰 수: 70,  고유어 수: 57
\end{lstlisting}
\end{frame}
